\section{An ordinary-angle C1 family: all independent tests}
\label{sec:c1-ordinary-angle}

This section closes an actual source family outside the small-$k$
uniform theorem. All inverses below are either real constants or
bounded units in one of the two independent real closed fields.
In particular no quotient of a testing infinitesimal by a source
infinitesimal is used.

\begin{theorem}[Regular ordinary-angle C1 source]
\label{thm:c1-ordinary-angle}
Use the seven native C1 coquotients and original coefficients of
the source-definition appendix, with $c=2$ and $\lambda=n=0$.
Suppose
\[
 k,s,a,t,M,J\text{ are positive bounded units},\qquad
 \eta>0\text{ is infinitesimal},\qquad
 f^2=\eta,\quad g^2=s^2\eta,
\]
\[
 K=1+k^2,\quad R=K-ak^2\eta,\quad T=K-k^2s^2\eta t,
 \quad M=R-t>0,\quad J=T-a>0.                 \tag{OA1}
\]
Then the full original parent is tensor SOS at every independent
projective test, for every real or complex state.
\end{theorem}

Here $s=g/f$ is the actual contact state ratio. Its being a unit is
part of the hypothesis. The source normalization and all six original
retained event coefficients are preserved. In particular this theorem
includes $k=s=1$, $a=t=1/2$, and arbitrary positive infinitesimal $\eta$.

\paragraph{Exact source and its leading scales.}
Put $e=k/s$ and $\Delta=s/k$ in the native formulas. To fix completely
which source is meant, they are
\[
\begin{gathered}
 d=\frac{2a\eta c_d}{c_3},\quad
 m=\frac{2ak^2c_2}{s^2Mc_3},\quad
 A=\frac{2a\eta^2c_4}{c_3},\quad
 \alpha=\frac{2as^2k^4\eta^2c_\alpha}{c_3},\\
 D+\frac d2=\frac{2ac_5}{Jc_3},\qquad
 \beta+\frac d2=\frac{2ak^2c_b}{s^2tc_3},\qquad c=2.
                                                        \tag{OA2}
\end{gathered}
\]
The letter $D$ here denotes the fifth event, whereas $D_u=1+u^2$.
Setting $\eta=\zeta=0$ in the short native polynomials gives
$c_2=c_b=c_3=c_5=c_4=c_\alpha=2K$ and $c_d=4K^2$.
Consequently the following positive bounded source expressions have
the same residues as the indicated original coefficients:
\[
\begin{gathered}
 m_0=\frac{2ak^2}{s^2(K-t)},\quad
 c_0=2,\quad D_0=\frac{2a}{K-a},\quad
 \beta_0=\frac{2ak^2}{s^2t},\\
 d/\eta\sim4aK,\qquad A/\eta^2\sim2a,\qquad
 \alpha/\eta^2\sim2as^2k^4.                         \tag{OA3}
\end{gathered}
\]
The symbol $\sim$ here means equality of residues, not equality of
the source. Each displayed inverse is a pure source unit inverse.
All original events are positive under (OA1), as follows from their
positive leading coefficients.

The exact contact is
\[
 u_c=k,\quad v_c=ks\eta,\quad
 (W_c,X_c,Y_c,Z_c)=(s,-k(1-ts^2\eta),ks(1-a\eta),1).
                                                        \tag{OA4}
\]
Its value and all first derivatives vanish. This assertion uses the
native identities (SD9)--(SD11), which are polynomial identities valid
for arbitrary $k$, not an asymptotic extension of a small-$k$ assertion.
The accompanying small checker repeats all six native balances.
For actual canonical C1 contacts, the source transport and nonzero
selected contact minor in (CE1)--(CE2) identify (OA2) with the original
source; no new source-entry assumption is implicit here.

\paragraph{A four-square reduction.}
Write
\[
 x=X+uZ,\quad y=uW-Y,\quad D_u=1+u^2,\quad D_v=1+v^2.
\]
The exact complete parent, including every signed term, is
\begin{align*}
 F={}&4m\{v(D_uW-uy)-x\}^2+4cD_vy^2
       +4AD_uD_v(uW-y)^2\\
 &+(4D+2d)\{y+v(ux-D_uZ)\}^2
       +4\alpha D_u\{v(uW-y)-Z\}^2\\
 &+(4\beta+2d)D_vx^2
       +2dD_u\bigl(\{v(uW-y)-Z\}^2-D_vZ^2\bigr).
                                                        \tag{OA5}
\end{align*}
The last line includes the essential $2dD_u(vY-Z)^2$ term.
No event is capped or discarded.

For a moment use formal variables $\xi,\chi,\nu$ and substitute
$x=\eta\xi$, $y=\eta\chi$, $v=\eta\nu$ into (OA5).
The coefficient of $\eta^2$ is
\begin{align*}
 \mathcal F={}&4m_0(\nu D_uW-\xi)^2+4\beta_0\xi^2
       +4c_0\chi^2+4D_0(\chi-\nu D_uZ)^2\\
 &+8aD_uu^2W^2+8as^2k^4D_uZ^2
       -16aKD_uu\nu WZ.
                                                        \tag{OA6}
\end{align*}
This is only a polynomial way to read a bounded Gram residue; it does
not substitute or compare any actual testing and coefficient rates.
Completing the first four squares gives
\[
 \mathcal F=4(m_0+\beta_0)
  \left(\xi-\frac tK\nu D_uW\right)^2
 +4(c_0+D_0)\left(\chi-\frac aK\nu D_uZ\right)^2
 +8aD_uS,                                               \tag{OA7}
\]
where
\[
 S=u^2W^2+s^2k^4Z^2+
   \frac{D_u}{K}\nu^2\left(\frac{k^2}{s^2}W^2+Z^2\right)
   -2Ku\nu WZ.                                         \tag{OA8}
\]

For bounded $u,v$ with $v$ infinitesimal, (OA5) is an exact bounded
quadratic Gram on
\[
                 x,\ y,\ vW,\ vZ,\ \eta W,\ \eta Z.    \tag{OA9}
\]
To see the assertion directly, expand only the displayed squares in
(OA5). Each monomial contains a product of two listed roots; all
remaining factors are bounded. Terms omitted in (OA6), after this
allocation, carry $v$, $\eta$, or a source coefficient difference of
zero residue. Hence its Gram residue is exactly (OA6)--(OA8), with
$\eta\nu$ read as $v$. This observation also applies to the modified
endpoint root lists below.

\paragraph{Angles away from the actual contact.}
Sign gauges allow $u,v\geq0$. Suppose first that $u$ has positive
finite residue. In the four roots
$\nu W,\nu Z,W,Z$, split the cross term in (OA8) between the pairs
$(\nu W,Z)$ and $(\nu Z,W)$, using their identical products.
The two positive cross-term capacities are
\[
 \sqrt{D_u/K}\,k^3,\qquad \sqrt{D_u/K}\,u.
\]
Their sum is strictly larger than $Ku$ unless $\bar u=\bar k$, since
the exact factorization is
\[
 (1+u^2)(u+k^3)^2-K^3u^2
  =(u-k)^2(k^4+2k^3u+2ku+u^2).                         \tag{OA10}
\]
Thus, when $\bar u\ne\bar k$, choose a real splitting constant
strictly within the two residue capacities. The resulting two
two-by-two blocks, together with the two completed hard squares,
give a strict residue Gram on (OA9).

If $u$ is infinitesimal, replace $\eta W$ in (OA9) by $\eta uW$.
Assign the entire cross term to $(vW,\eta Z)$; its coefficient has
zero residue. The four soft diagonal coefficients are positive units.
The same monomial allocation in (OA5) proves an exact strict Gram.

The reciprocal Alice endpoint requires no source reflection. Put
$h=1/u$, multiply the homogeneous parent by the pure testing square
$h^2$, and use $x'=hX+Z$, $y'=W-hY$. Its six roots are
\[
                 x',\ y',\ vY,\ vX,\ \eta Y,\ \eta hX.
                                                        \tag{OA11}
\]
When $h$ is infinitesimal, the four strong residue terms are
$4m_0(vY-x')^2+4\beta_0x'^2+4c_0y'^2+4D_0(y'+vX)^2$.
The two weak diagonal reserves are $8a(\eta Y)^2$ and
$8as^2k^4(\eta hX)^2$. The signed cross term is assigned as a
multiple of $h(vX)(\eta Y)$, so has zero residue; all other corrections
also have zero residue. This is a strict six-root Gram, including
the literal endpoint $h=0$ by the homogeneous formula.

\paragraph{Matched angle but unmatched state.}
It remains to consider $\bar u=\bar k$. At $u=k$, (OA8) factors as
\[
 S=(\nu Z-kW)^2+
       \left(\frac{k}{s}\nu W-sk^2Z\right)^2.          \tag{OA12}
\]
Use the exact roots
\[
 r_1=vZ-\eta kW,\qquad
 r_2=\frac{k}{s}vW-\eta sk^2Z,
\]
together with the two completed hard roots
$h_1=x-(t/K)D_uvW$, $h_2=y-(a/K)D_uvZ$, and $\eta$.
All six roots in (OA9) are bounded linear combinations of these five.
The exact relation
\[
 Wr_1-\frac{s}{k}Zr_2
   =\eta k(s^2Z^2-W^2)=p\eta                         \tag{OA13}
\]
shows that, if $\bar p\ne0$, adding the identically zero square
$(p\eta-Wr_1+(s/k)Zr_2)^2$ makes the five-root Gram strict.
This operation never inverts $p$.

After pure testing state normalization, any state with positive
parent residue already has a one-root certificate. Otherwise the
four strong events force $\bar x=\bar y=0$, so $W$ or $Z$ is a unit.
If $\bar p=0$, then $Z$ is a unit and we may put $Z=1$.
The remaining choices are $\bar W=\pm\bar s$. The negative choice
is a direct nonnegative-core exit. Indeed set
$w=(W+uY)/D_u=W-uy/D_u$. The pure testing core Schur calculation has hard block
\[
 \begin{pmatrix}1+v^2+u^2v^2&uv\\uv&(D_u+v^2)/D_u\end{pmatrix}\succ0
\]
and positive Schur denominator
$\mathcal D=u^2v^4+u^2v^2+u^2+v^4+2v^2+1$.
Its remaining scalar has positive factor $uvD_uD_v/\mathcal D$
and is proportional to
\[
 D_uuv w^2-2D_uD_vw+uv.
\]
For $u,v>0$ it is nonnegative if $w\leq0$. Hence a negative-core
test must have $w>0$; with $\bar y=0$ it cannot have
$\bar W=-\bar s$. A nonnegative core is itself a pure testing
square after evaluation, and each original event is already a square.
Literal $u=0$ or $v=0$ is included in this direct exit.

\paragraph{The contact: five roots and a finite Taylor argument.}
We have reduced to $Z=1$ and
\[
 u-k,\ W-s,\ v,\ x,\ y\text{ infinitesimal}.
\]
Set
\[
 x_c=\eta kts^2,\quad y_c=\eta ksa,\quad v_c=\eta ks,
\quad U=u-k,\quad B=W-s,
\]
\[
 N_1=x-x_c,\quad N_2=y-y_c,\quad N_3=v-v_c.
\]
The exact roots for this chart are
\[
                 N_1,\ N_2,\ N_3,\ \eta U,\ \eta B.   \tag{OA14}
\]
Here is the full allocation argument, including the normal terms.
As a polynomial in $x,y,v$, (OA5) belongs to
\[
              (x,y,v)^2+\eta(x,y,v)+(\eta^2)
\]
over the bounded polynomial coefficient ring in $u,W$ and the source
unit inverses already listed. Since $x_c,y_c,v_c$ are $\eta$ times
bounded source elements, the shifted polynomial belongs to
$(N_1,N_2,N_3)^2+\eta(N_1,N_2,N_3)+(\eta^2)$.
Its constant normal coefficient is $\eta^2$ times a polynomial in
$U,B$ whose constant and linear terms vanish by exact contact and
angle stationarity. Its linear normal coefficients are $\eta$ times
polynomials in $U,B$ whose constant terms vanish by exact state and
Bob stationarity. These conclusions can be read in the coefficient
field, where $\eta\ne0$; the resulting quotients are bounded by the
displayed factorization, not by cancellation of a mixed element.
Monic Taylor division by $U,B$ therefore allocates these terms to
products of the roots (OA14). Terms of normal degree at least two
are allocated to their first two normal factors. This is an exact
finite Gram construction with bounded entries.

All terms beyond weighted degree two carry one of the infinitesimal
deviations $N_i,U,B$, or a coefficient difference of zero residue.
The Gram residue is consequently the weighted Hessian of (OA6) at
\[
 (u,W,\nu,\xi,\chi)=(k,s,ks,kts^2,ksa),\qquad Z=1.
\]
Its $\xi,\chi$ block is positive definite by (OA7).
The Schur complement is $8aK$ times the half-Hessian of $S$ in
$(\nu,W,u)$ at this point. That half-Hessian is
\[
 \begin{pmatrix}
 K&k(k^2-1)&s(k^2-1)\\
 k(k^2-1)&k^2K&sk(k^4+1)/K\\
 s(k^2-1)&sk(k^4+1)/K&s^2K
 \end{pmatrix},                                      \tag{OA15}
\]
with leading principal minors
\[
                    K,\qquad4k^4,\qquad12s^2k^6/K.
                                                        \tag{OA16}
\]
These are positive units. The five-root Gram is strict, which proves
the contact chart without assuming that $v/\eta$, $x/\eta$, or
$y/\eta$ is bounded.

\paragraph{Exhaustion and restoration.}
For completeness, if the homogeneous Bob sine $v_1$ has nonzero
residue, the four strong packet amplitudes have trivial common state
kernel at every projective Alice angle. Indeed $x=y=0$ and the other
two amplitudes then give
$v_1(u_0W+u_1Y)=v_1(u_1X-u_0Z)=0$; since
$u_0^2+u_1^2$ is a testing unit, all state coordinates vanish.
Their four-root Gram absorbs the infinitesimal signed coefficient;
all other event squares may be retained. Thus only the Bob-small
chart used above remains. The finite and reciprocal Alice charts,
the two state branches, and the exact contact chart exhaust it.
Every strict bounded residue Gram is tensor SOS by the established
strict-Gram lemma. All projective multipliers are pure testing squares
and are restored by separate field square roots. Complex states split
into their real and imaginary parts.

\paragraph{Exact advance and remaining scope.}
The theorem closes the regular ordinary-$k$, comparable-$f,g$ source
slice (OA1), including the previous explicit $k=1$ entry obstruction.
It supplies a general mechanism: two completed strong squares, a
scalar capacity factorization locating the unique angular degeneracy,
and finite weighted contact division with a strict five-root Hessian.
It makes no claim for a nonunit $s$, a nonunit shape or gap among
$a,t,M,J$, non-small $\eta$ or $\zeta$, or the still-unclosed
small-$k$ unit/unbounded-$\eta$ source rows.
